% GENERATED FILE. Edit canonical lessons, then run npm run generate:books.
% canonical-source-hash: fnv1a64:6a247f3e92236481
% canonical-lessons: PT-C418-bote, PT-C418-remo, PT-C418-remar, PT-C418-canoa, PT-C418-barqueiro

\chapter{De barco}
\label{ch:pt-de-barco}
\hlchaptermodality{418}

\begin{chapteropening}
\textbf{By the end of this chapter:} \emph{I can talk about small boats: the dinghy, the oar, rowing, the canoe and the boatman.}

Complete the last lesson of chapter 418: I can talk about small boats: the dinghy, the oar, rowing, the canoe and the boatman.
\end{chapteropening}

\section[o bote]{o bote — the dinghy, the small boat}

\label{lesson:PT-C418-bote}



\begin{quote}
Before the new one: say the Portuguese for to go ashore, then the Portuguese for to set sail.
\end{quote}

\subsection*{You'll want to know: o bote}
\textbf{o bote} — \textquotedblleft{}the dinghy, the small boat\textquotedblright{}.

\begin{quote}
Fomos de bote até à praia do outro lado.
\end{quote}

\begin{quote}
\emph{We went by dinghy to the beach on the other side.}
\end{quote}

\emph{O bote salva-vidas} — the lifeboat.

\begin{cousinweb}
\emph{bote} came from Old French \emph{bot}, which was taken from Middle English \emph{boot}, Old English \emph{bāt} 'boat' (EtymDB).
\end{cousinweb}

\begin{grammarlens}[title={What you've built}]
One more word for telling, explaining and describing at B1.
\end{grammarlens}

\subsection*{Your turn}
\begin{itemize}
\raggedright
  \item \emph{Say it:} \emph{o bote}
  \item \emph{Say it:} \emph{o bote}, once more
  \item \emph{Say it:} a sentence of your own with \emph{o bote}
  \item \emph{Recall:} say the Portuguese for to go ashore, then the Portuguese for to set sail, then say \emph{o bote} again
\end{itemize}

\begin{tcolorbox}[breakable,colback=teal!4,colframe=teal!35!black,title={Before you move on}]
What does o bote mean? (\textquotedblleft{}The dinghy, the small boat\textquotedblright{}.) Say it once more.
\end{tcolorbox}
\section[o remo]{o remo — the oar}

\label{lesson:PT-C418-remo}



\begin{quote}
Before the new one: say the Portuguese for to set sail, then the Portuguese for the dinghy.
\end{quote}

\subsection*{You'll want to know: o remo}
\textbf{o remo} — \textquotedblleft{}the oar\textquotedblright{}.

\begin{quote}
Um dos remos caiu à água.
\end{quote}

\begin{quote}
\emph{One of the oars fell into the water.}
\end{quote}

In sport: \emph{o remo} — rowing.

\begin{cousinweb}
\emph{remo} is inherited from Latin \emph{rēmus} 'oar' (EtymDB).
\end{cousinweb}

\begin{grammarlens}[title={What you've built}]
One more word for telling, explaining and describing at B1.
\end{grammarlens}

\subsection*{Your turn}
\begin{itemize}
\raggedright
  \item \emph{Say it:} \emph{o remo}
  \item \emph{Say it:} \emph{o remo}, once more
  \item \emph{Say it:} a sentence of your own with \emph{o remo}
  \item \emph{Recall:} say the Portuguese for to set sail, then the Portuguese for the dinghy, then say \emph{o remo} again
\end{itemize}

\begin{tcolorbox}[breakable,colback=teal!4,colframe=teal!35!black,title={Before you move on}]
What does o remo mean? (\textquotedblleft{}The oar\textquotedblright{}.) Say it once more.
\end{tcolorbox}
\section[remar]{remar — to row}

\label{lesson:PT-C418-remar}



\begin{quote}
Before the new one: say the Portuguese for the dinghy, then the Portuguese for the oar.
\end{quote}

\subsection*{You'll want to know: remar}
\textbf{remar} — \textquotedblleft{}to row\textquotedblright{}.

\begin{quote}
Remámos durante uma hora no lago.
\end{quote}

\begin{quote}
\emph{We rowed for an hour on the lake.}
\end{quote}

\emph{O remador} — the rower.

\begin{cousinweb}
\emph{remar} is built on \emph{remo} 'oar', which is inherited from Latin \emph{rēmus} (EtymDB).
\end{cousinweb}

\begin{grammarlens}[title={What you've built}]
One more word for telling, explaining and describing at B1.
\end{grammarlens}

\subsection*{Your turn}
\begin{itemize}
\raggedright
  \item \emph{Say it:} \emph{remar}
  \item \emph{Say it:} \emph{remar}, once more
  \item \emph{Say it:} a sentence of your own with \emph{remar}
  \item \emph{Recall:} say the Portuguese for the dinghy, then the Portuguese for the oar, then say \emph{remar} again
\end{itemize}

\begin{tcolorbox}[breakable,colback=teal!4,colframe=teal!35!black,title={Before you move on}]
What does remar mean? (\textquotedblleft{}To row\textquotedblright{}.) Say it once more.
\end{tcolorbox}
\section[a canoa]{a canoa — the canoe}

\label{lesson:PT-C418-canoa}



\begin{quote}
Before the new one: say the Portuguese for the oar, then the Portuguese for to row.
\end{quote}

\subsection*{You'll want to know: a canoa}
\textbf{a canoa} — \textquotedblleft{}the canoe\textquotedblright{}.

\begin{quote}
Descemos o rio de canoa.
\end{quote}

\begin{quote}
\emph{We went down the river by canoe.}
\end{quote}

A light, narrow boat that you paddle.

\begin{cousinweb}
\emph{canoa} was borrowed through Spanish from Taíno \emph{kanoa} 'dugout canoe' (EtymDB), a word of the Caribbean islands.
\end{cousinweb}

\begin{grammarlens}[title={What you've built}]
One more word for telling, explaining and describing at B1.
\end{grammarlens}

\subsection*{Your turn}
\begin{itemize}
\raggedright
  \item \emph{Say it:} \emph{a canoa}
  \item \emph{Say it:} \emph{a canoa}, once more
  \item \emph{Say it:} a sentence of your own with \emph{a canoa}
  \item \emph{Recall:} say the Portuguese for the oar, then the Portuguese for to row, then say \emph{a canoa} again
\end{itemize}

\begin{tcolorbox}[breakable,colback=teal!4,colframe=teal!35!black,title={Before you move on}]
What does a canoa mean? (\textquotedblleft{}The canoe\textquotedblright{}.) Say it once more.
\end{tcolorbox}
\section[o barqueiro]{o barqueiro — the boatman, the ferryman}

\label{lesson:PT-C418-barqueiro}



\begin{quote}
Before the new one: say the Portuguese for to row, then the Portuguese for the canoe.
\end{quote}

\subsection*{You'll want to know: o barqueiro}
\textbf{o barqueiro} — \textquotedblleft{}the boatman, the ferryman\textquotedblright{}.

\begin{quote}
O barqueiro leva as pessoas de uma margem à outra.
\end{quote}

\begin{quote}
\emph{The boatman takes people from one bank to the other.}
\end{quote}

Built on \emph{o barco}.

\begin{cousinweb}
\emph{barqueiro} is built on \emph{barco} 'boat' with \emph{-eiro}, a person who works with it; the family goes back to Late Latin \emph{barca} (Lewis \& Short).
\end{cousinweb}

\begin{grammarlens}[title={What you've built}]
One more word for telling, explaining and describing at B1.
\end{grammarlens}

\subsection*{Your turn}
\begin{itemize}
\raggedright
  \item \emph{Say it:} \emph{o barqueiro}
  \item \emph{Say it:} \emph{o barqueiro}, once more
  \item \emph{Say it:} a sentence of your own with \emph{o barqueiro}
  \item \emph{Recall:} say the Portuguese for to row, then the Portuguese for the canoe, then say \emph{o barqueiro} again
\end{itemize}

\begin{tcolorbox}[breakable,colback=teal!4,colframe=teal!35!black,title={Before you move on}]
What does o barqueiro mean? (\textquotedblleft{}The boatman, the ferryman\textquotedblright{}.) Say it once more.
\end{tcolorbox}
